\section{Cross-Language Replication: WMT14 en-de}
\label{sec:ende}

This section asks whether the en-fr 2$\times$2 pattern generalizes to a different language pair. We pick WMT14 en-de --- a $4.5\times$ smaller raw corpus than en-fr (4.5M vs.\ 40.8M), where the raw data is also already cleaner per-pair. The corpus is small enough that we cannot construct a clean v1/v2 split (there is no ``loose'' option that yields meaningfully more data); instead we run a single Base + a single Big on the standard cleaned corpus and look at the Base/Big gap.

\subsection{Corpus and setup}
WMT14 en-de download via HuggingFace Datasets (4{,}509{,}489 raw pairs); cleaned with the v1 thresholds; 4{,}174{,}104 pairs survive (92.6\% keep rate). SPM 32K joint vocabulary, \texttt{character\_coverage=1.0} (preserves German umlauts). Train/valid/test = 4.17M / 3{,}000 / 3{,}003.

Hyperparameters mirror the en-fr Base config (\S\ref{sec:setup}) with two changes: max\_steps from 100K $\rightarrow$ 230K (extended after the curve at 150K showed continued non-trivial improvement), and dropout 0.1 $\rightarrow$ 0.3 for Big (following the original Vaswani recommendation for the smaller en-de corpus).

\subsection{Base en-de}
Base en-de trains for 230K steps ($\sim$~6 epochs at this batch size), best valid 24.25 (single ckpt), 5-checkpoint averaged valid 24.35 / test 24.04 with default beam=5 lp=1.0. The original paper reports tokenized BLEU 27.3 on en-de Base; converting via the ${\sim}1.5$--$2$ tokenized-vs-sacrebleu band, our 24.04 sacrebleu is $\sim$~25.5--26 tokenized BLEU --- 1--1.5 BLEU below the original report. Following the original paper's en-de decoding hyperparameters (beam=4, length-penalty=0.6) does not close this gap (24.19 sacrebleu, +0.15). The gap is small relative to the +0.56 / -0.87 effects we make claims on, and likely reflects differences in BPE tokenization details and training-corpus filtering that we cannot precisely reproduce.

\subsection{Big en-de}
Big en-de trains for 466K steps with dropout 0.3 (original Vaswani's en-de regularization), best single-checkpoint valid 23.41, 5-checkpoint averaged valid 23.01 / test 22.47.

\subsection{The pattern replicates}
\begin{table}[t]
\centering
\small
\begin{tabular}{lrrr}
\toprule
\textbf{en-de} & \textbf{Base} & \textbf{Big} & \textbf{$\Delta$} \\
\midrule
Valid & 24.35 & 23.01 & $-1.34$ \\
Test  & 24.04 & 22.47 & $-1.57$ \\
\bottomrule
\end{tabular}
\caption{en-de results: Big falls below Base by 1.5 BLEU on test newstest2014. The pattern qualitatively matches Big v2's loss against Base v2 on en-fr noisy data ($-0.87$); we read both as the small-clean-data regime where 4.5M is too small for Big's 209M parameters to find a winning regime.}
\label{tab:ende}
\end{table}

The Base $>$ Big ordering is consistent with the en-fr v2 result: at 4.5M effective pairs, the data is small enough that Big over-parameterizes relative to it, and the resulting capacity overhang costs BLEU. This is a different mechanism from v2's noise-amplification --- 4.5M en-de is genuinely cleaner than 30M en-fr v2 --- but it leads to the same Big$<$Base ordering.

The 4.5M scale is roughly half of v1.1's 9.3M, and unlike v1.1, it is on the wrong side of where Big becomes useful. We do not claim a precise scale threshold from a single en-de data point; we do claim that the en-de result is consistent with a unifying picture: Big needs (a) clean data, (b) enough of it. v1.1 has both; v2 has (b) but not (a); en-de has (a) but not (b). All three exhibit Base $\geq$ Big.

\subsection{HuggingFace release}
Both en-de checkpoints are published on the HuggingFace Hub:
\texttt{euswbnix/transformer-wmt14-ende-base} (60M, test 24.04) and
\texttt{euswbnix/transformer-wmt14-ende-big} (209M, test 22.47).
